% Scanning Compressed Strings: DFA-Based LIKE Pushdown for FSST-Encoded Columnar Data
% Target: DaMoN (Data Management on New Hardware) @ SIGMOD
% Format: 6-8 pages, double column, ACM sigconf

\documentclass[sigconf,nonacm]{acmart}

\usepackage{booktabs}
\usepackage{algorithm}
\usepackage{algpseudocode}
\usepackage{subcaption}
\usepackage{xspace}

% Convenience macros
\newcommand{\fsst}{\textsc{fsst}\xspace}
\newcommand{\like}{\texttt{LIKE}\xspace}
\newcommand{\vortex}{\textsc{Vortex}\xspace}
\newcommand{\dfa}{\textsc{dfa}\xspace}
\newcommand{\kmp}{\textsc{kmp}\xspace}

\begin{document}

\title{Scanning Compressed Strings: DFA-Based LIKE Pushdown\\for FSST-Encoded Columnar Data}

\author{TODO}
\affiliation{%
  \institution{TODO}
}
\email{TODO}

% ===========================================================================
% ABSTRACT
% ===========================================================================

\begin{abstract}
Modern columnar stores compress strings with FSST (Fast Static Symbol Table),
which replaces recurring byte sequences with single-byte codes using a learned
255-entry symbol table. Evaluating predicates like \texttt{LIKE~'prefix\%'} or
\texttt{LIKE~'\%needle\%'} conventionally requires decompressing the data first,
paying the cost of memory allocation, cache pollution, and memcpy before the
actual comparison can begin.

We present a technique that evaluates LIKE predicates directly on
FSST-compressed code streams by constructing a compact DFA over the symbol
table. The core challenge is that FSST symbols are variable-length (1--8 bytes),
so a single code byte may advance the matcher through multiple byte-level
states. Our construction addresses this in four stages: (1)~build a KMP failure
function on the needle bytes, (2)~materialize a byte-level transition table,
(3)~simulate each symbol's multi-byte expansion to produce symbol-level
transitions---the step with no analogue in standard DFA literature---and
(4)~flatten the result into a cache-resident \texttt{u8}-indexed lookup table.
Escape codes (literal bytes without a symbol) are handled via a sentinel
architecture that keeps the hot path branchless while gracefully degrading as
escape rate increases.

The resulting DFA fits entirely in L1 cache ($\leq$64\,KB for patterns up to
253 bytes), requires zero heap allocation during scanning, and processes one
FSST code per iteration rather than one byte. On real-world datasets from
ClickBench and FineWeb, FSST achieves 3--5$\times$ compression with mean symbol
lengths of 3--5 bytes and escape rates below 5\%. For \texttt{prefix\%}
patterns the DFA matches or exceeds Arrow's SIMD string kernel while skipping
decompression entirely; for \texttt{\%needle\%} patterns it delivers consistent
speedups that scale with compression ratio. We implement the technique in
Vortex, an open-source compressed columnar engine, where it composes with
zone-map pruning: the zone map eliminates non-matching row groups, and the DFA
scans survivors without materializing strings.
\end{abstract}

\maketitle

% ===========================================================================
% 1. INTRODUCTION
% ===========================================================================

\section{Introduction}
\label{sec:intro}

% WHY: String-heavy analytical queries are increasingly common, and FSST is
% the dominant lightweight string compression in columnar stores.

% PROBLEM: LIKE predicates require decompression today — this is wasteful
% because the compressed representation already contains enough information
% to evaluate the predicate.

% CONTRIBUTION: We show how to construct a DFA that operates directly on
% FSST code streams, and demonstrate significant speedups on real data.

% OUTLINE: Brief roadmap of the paper.

\textbf{TODO.}

% ===========================================================================
% 2. BACKGROUND
% ===========================================================================

\section{Background}
\label{sec:background}

% 2.1 FSST Encoding
% - Symbol table: up to 255 symbols, each 1-8 bytes
% - Compression: greedy left-to-right, longest-match
% - Escape code (255): literal byte follows
% - Typical compression ratios on string data: 3-5x

\subsection{FSST Encoding}
\label{sec:bg-fsst}

\textbf{TODO.}

% 2.2 LIKE Predicate Semantics
% - SQL LIKE with % and _ wildcards
% - We focus on two common cases:
%   (a) prefix%  — starts-with
%   (b) %needle% — contains
% - These cover the majority of LIKE usage in practice

\subsection{LIKE Predicate Semantics}
\label{sec:bg-like}

\textbf{TODO.}

% 2.3 KMP String Matching
% - Failure function, linear-time matching
% - Why KMP (not Boyer-Moore): we process symbols left-to-right,
%   matching FSST's encoding direction

\subsection{KMP String Matching}
\label{sec:bg-kmp}

\textbf{TODO.}

% ===========================================================================
% 3. DFA CONSTRUCTION
% ===========================================================================

\section{DFA Construction}
\label{sec:construction}

% Overview: four-stage pipeline that transforms a LIKE pattern + FSST symbol
% table into a flat lookup table. Walk through each stage with a running example.

% Running example:
%   needle = "abab"
%   symbols: s0="ab" (code 0), s1="ac" (code 1), s2="b" (code 2)

\textbf{TODO: introduce running example.}

% --------------------------------------------------------------------------
% Stage 1: KMP failure function
% --------------------------------------------------------------------------

\subsection{Stage 1: KMP Failure Function}
\label{sec:stage1}

% Build standard KMP failure function on needle bytes.
% For needle "abab": failure = [0, 0, 1, 2]
% This enables O(n) matching with backtracking on mismatch.

\textbf{TODO.}

% --------------------------------------------------------------------------
% Stage 2: Byte-level transition table
% --------------------------------------------------------------------------

\subsection{Stage 2: Byte-Level Transition Table}
\label{sec:stage2}

% Materialize (state × byte) → state table.
% States: 0..needle_len (progress states) + ACCEPT + possibly FAIL (prefix only).
% For prefix%: mismatch → sticky FAIL (no backtracking).
% For %needle%: mismatch → KMP failure fallback.

\textbf{TODO.}

% --------------------------------------------------------------------------
% Stage 3: Symbol-level transitions (CORE NOVELTY)
% --------------------------------------------------------------------------

\subsection{Stage 3: Symbol-Level Transitions}
\label{sec:stage3}

% THE KEY STEP: For each (state, symbol_code) pair, feed the symbol's bytes
% through the byte-level table and record the final state.
%
% This handles:
%   - Multi-byte symbols advancing through multiple match states
%   - Cross-symbol boundary matching (pattern straddles two symbols)
%   - KMP fallback within a single symbol (e.g., "abac" symbol vs "abab" needle)
%
% Worked example with the running example symbols.
%
% Correctness argument: the simulation produces identical results to
% byte-by-byte matching on decompressed data, regardless of how the
% compressor chose to decompose the string into symbols.

\textbf{TODO.}

% --------------------------------------------------------------------------
% Stage 4: Flat u8 table with escape sentinel
% --------------------------------------------------------------------------

\subsection{Stage 4: Flat Table with Escape Sentinel}
\label{sec:stage4}

% Flatten into table[state * 256 + code] → next_state.
% State IDs are u8 → table ≤ 64KB → fits in L1.
%
% Escape handling:
%   - Code 255 (FSST escape) maps to sentinel value in the symbol table
%   - Scanner detects sentinel, reads next literal byte
%   - Looks up in separate byte-level escape table
%   - Hot path (codes 0-254): one branchless table lookup
%   - Cold path (code 255): one branch + one extra lookup
%
% Pattern length limits:
%   - prefix%: up to 253 bytes (N progress + accept + fail + sentinel ≤ 256)
%   - %needle%: up to 254 bytes (N progress + accept + sentinel ≤ 256)

\textbf{TODO.}

% --------------------------------------------------------------------------
% Pseudocode
% --------------------------------------------------------------------------

\subsection{Scan Algorithm}
\label{sec:scan}

% Present the scanning algorithm as pseudocode.
% Emphasize the hot path: while (pos < len) { state = table[state][codes[pos++]]; }
% Show how escape sentinel interrupts this.

\begin{algorithm}
\caption{DFA scan over FSST-compressed codes}\label{alg:scan}
\begin{algorithmic}[1]
\Require{code stream $C[0..n)$, symbol table $T$, escape table $E$, sentinel $\sigma$}
\Ensure{true if pattern matches}
\State $s \gets s_0$ \Comment{start state}
\State $i \gets 0$
\While{$i < n$}
    \State $s' \gets T[s][C[i]]$
    \If{$s' = \sigma$} \Comment{escape: cold path}
        \State $i \gets i + 1$
        \State $s \gets E[s][C[i]]$
    \Else \Comment{symbol: hot path}
        \State $s \gets s'$
    \EndIf
    \If{$s = s_{\text{accept}}$} \Return true
    \EndIf
    \State $i \gets i + 1$
\EndWhile
\State \Return false
\end{algorithmic}
\end{algorithm}

% ===========================================================================
% 4. INTEGRATION
% ===========================================================================

\section{Integration with Vortex}
\label{sec:integration}

% 4.1 Pushdown framework
% - Vortex evaluates scalar functions on encoded arrays without decoding
% - The LIKE kernel checks if the array is FSST-encoded, attempts DFA
%   construction, falls back to decompression if unsupported pattern

% 4.2 Zone-map composition
% - Vortex files store per-zone min/max statistics
% - Zone maps prune row groups that cannot match the LIKE pattern
% - DFA scans only surviving zones
% - Combined benefit: zone maps reduce volume, DFA reduces cost per byte

\textbf{TODO.}

% ===========================================================================
% 5. EVALUATION
% ===========================================================================

\section{Evaluation}
\label{sec:eval}

% Setup: machine specs, compiler, datasets, methodology (samples × iters)

\subsection{Experimental Setup}
\label{sec:eval-setup}

\textbf{TODO: machine specs, compiler version, methodology.}

% --------------------------------------------------------------------------
% 5.1 Symbol table characterization (Experiment A)
% --------------------------------------------------------------------------

\subsection{Symbol Table Characterization}
\label{sec:eval-symtab}

% Real-world datasets: ClickBench (URL, Title, SearchPhrase, Referer),
% FineWeb (url, text), TPC-H (l_comment).
%
% Table: n_strings, avg_len, n_symbols, mean_sym_len, compression_ratio,
% escape_rate, entropy, effective_alphabet, p50/p90 symbol coverage.
%
% Key finding: escape rates are consistently low (<5%) on real data,
% validating the sentinel architecture.

\textbf{TODO: summary table of symbol table characteristics.}

% --------------------------------------------------------------------------
% 5.2 DFA vs decompress-then-match (Experiment B)
% --------------------------------------------------------------------------

\subsection{DFA vs Decompress-then-Match}
\label{sec:eval-throughput}

% Bar chart: speedup per dataset × pattern type (prefix%, %needle%).
% Use real patterns on real data.
% Show Arrow LIKE kernel as the baseline (not a strawman).
%
% Key finding: DFA wins most convincingly on well-compressed data with
% longer symbols (URLs). On high-entropy data (SearchPhrase), the advantage
% is smaller but still positive.

\textbf{TODO: throughput comparison, bar chart.}

% --------------------------------------------------------------------------
% 5.3 Escape rate sensitivity (Experiment C)
% --------------------------------------------------------------------------

\subsection{Escape Rate Sensitivity}
\label{sec:eval-escape}

% Line chart: DFA throughput (GB/s) vs escape rate (0-50%).
% Controlled sweep using noise injection on FineWeb text.
%
% Key finding: graceful degradation — throughput drops linearly, no cliff.
% Even at 20% escape rate, DFA still competitive with decompression.

\textbf{TODO: escape rate sweep, line chart.}

% --------------------------------------------------------------------------
% 5.4 Symbol length vs throughput (Experiment D)
% --------------------------------------------------------------------------

\subsection{Mean Symbol Length vs Throughput}
\label{sec:eval-symlen}

% Scatter: mean_sym_len (x) vs throughput (y) for DFA and decompress.
% Real datasets as named points overlaid.
%
% Key finding: DFA throughput scales roughly linearly with mean symbol length
% (fewer transitions per string). Decompression throughput is mostly flat.

\textbf{TODO: symbol length scatter plot.}

% --------------------------------------------------------------------------
% 5.5 Pattern length scaling (Experiment E)
% --------------------------------------------------------------------------

\subsection{Pattern Length Scaling}
\label{sec:eval-patlen}

% Line chart: throughput vs pattern length (1-200 bytes).
% DFA table size grows linearly; does it cause cache pressure?
%
% Key finding: throughput roughly constant until table exceeds L1 cache.

\textbf{TODO: pattern length sweep.}

% --------------------------------------------------------------------------
% 5.6 Construction cost (Experiment F)
% --------------------------------------------------------------------------

\subsection{Construction Cost}
\label{sec:eval-construction}

% Table: construction time (ns) for various pattern lengths and types.
% Compare against scan time for 1K, 10K, 100K rows.
%
% Key finding: construction is <10µs, amortized after a few hundred rows.

\textbf{TODO: construction cost table.}

% --------------------------------------------------------------------------
% 5.7 End-to-end with zone pruning (Experiment G)
% --------------------------------------------------------------------------

\subsection{End-to-End File Scan}
\label{sec:eval-e2e}

% Stacked bar: query time breakdown (zone prune + DFA scan + other).
% Full Vortex pipeline on ClickBench/FineWeb data written to Vortex files.
%
% Key finding: zone maps + DFA compose multiplicatively.

\textbf{TODO: end-to-end results.}

% --------------------------------------------------------------------------
% 5.8 Micro-architectural analysis (Experiment I)
% --------------------------------------------------------------------------

\subsection{Micro-Architectural Analysis}
\label{sec:eval-perf}

% perf counters: L1/L2/L3 cache misses, branch mispredictions, IPC.
% Compare DFA scan vs decompress-then-match.
%
% Key finding: DFA has dramatically fewer cache misses (small working set)
% and fewer branch mispredictions (branchless hot path).

\textbf{TODO: perf counter comparison.}

% ===========================================================================
% 6. RELATED WORK
% ===========================================================================

\section{Related Work}
\label{sec:related}

% 6.1 FSST
% - Boncz et al., VLDB 2020: the FSST paper
% - Used in DuckDB, Velox, Vortex, and others

% 6.2 Predicate evaluation on compressed data
% - Dictionary pushdown (trivial: compare against dictionary entries)
% - Run-length encoding pushdown
% - Bit-packed integer pushdown
% - None of these handle variable-length string compression like FSST

% 6.3 Compressed pattern matching
% - Navarro & Tarhio (2005): pattern matching on LZ77
% - Amir, Benson, Farach (1996): matching on LZW
% - These target different compression schemes and are not directly applicable

% 6.4 SIMD string matching
% - Arrow LIKE kernel uses SIMD for prefix/contains
% - Our DFA operates at a coarser granularity (codes vs bytes)
% - Complementary: could SIMD-ize the DFA table lookups

\textbf{TODO.}

% ===========================================================================
% 7. CONCLUSION
% ===========================================================================

\section{Conclusion}
\label{sec:conclusion}

% Summary of contributions:
% 1. Four-stage DFA construction for LIKE on FSST-compressed data
% 2. Sentinel escape architecture for branchless hot path
% 3. Comprehensive evaluation on real-world datasets
%
% Future work:
% - Extend to regex (fused DFA prototype already works)
% - SIMD-ize the DFA table lookups (gather instructions)
% - Extend to other LIKE patterns (suffix, interior wildcards)
% - Apply to other variable-length encodings beyond FSST

\textbf{TODO.}

% ===========================================================================
% REFERENCES
% ===========================================================================

\bibliographystyle{ACM-Reference-Format}
\bibliography{references}

\end{document}
